\documentclass[11pt,twocolumn]{article}

\usepackage[margin=1in]{geometry}
\usepackage{amsmath,amssymb,amsthm}
\usepackage{graphicx}
\usepackage{hyperref}
\usepackage{listings}
\usepackage{xcolor}
\usepackage{booktabs}
\usepackage{enumitem}
\usepackage{algorithm}
\usepackage{algpseudocode}

\hypersetup{colorlinks=true,linkcolor=blue!60!black,citecolor=green!40!black,urlcolor=blue!60!black}

\lstset{
  language=Rust,
  basicstyle=\ttfamily\scriptsize,
  keywordstyle=\color{blue!70!black}\bfseries,
  commentstyle=\color{green!50!black}\itshape,
  stringstyle=\color{red!60!black},
  breaklines=true,
  frame=single,
  backgroundcolor=\color{gray!5}
}

\newtheorem{definition}{Definition}
\newtheorem{theorem}{Theorem}
\newtheorem{proposition}{Proposition}

\title{Post-Quantum Privacy for Decentralized Governance:\\ML-KEM-768 Handshake and PQC-Encrypted Mail}
\author{Tristan Stoltz \\ Luminous Dynamics \\ \texttt{tristan.stoltz@evolvingresonantcocreationism.com}}
\date{March 2026}

\begin{document}

\twocolumn[
\begin{@twocolumnfalse}
\maketitle

\begin{abstract}
We present a post-quantum privacy architecture for decentralized governance
systems, implemented across two Mycelix components: the Identity cluster's
hybrid key infrastructure (Ed25519 + ML-KEM-768) and the Mail cluster's 12-zome
PQC-encrypted email system. The architecture provides three layers of quantum
resistance: (1) a hybrid handshake protocol combining classical Diffie-Hellman
with ML-KEM-768 key encapsulation, ensuring security even if one scheme is
broken; (2) AES-256-GCM message encryption using keys derived from the hybrid
handshake; and (3) forward secrecy through ephemeral key pairs per conversation.
We address the practical challenges of deploying PQC in a decentralized
context: larger key sizes (1,184 bytes for ML-KEM-768 vs. 32 bytes for Ed25519),
DHT bandwidth implications, and backward compatibility with agents that have
not yet upgraded to post-quantum keys. The system is integrated with
consciousness-gated access control, where encryption key exchange requires
Participant+ tier to prevent anonymous spam.
\end{abstract}
\end{@twocolumnfalse}
]

\section{Introduction}

The advent of cryptographically relevant quantum computers (CRQCs) threatens
the security foundations of all current decentralized systems. RSA and
elliptic-curve cryptography, which underpin blockchain signatures, TLS
connections, and identity verification, are vulnerable to Shor's
algorithm~\cite{shor1994}.

For governance systems, the quantum threat is particularly severe:

\begin{enumerate}[nosep]
  \item \textbf{Identity forgery}: An attacker with a CRQC could forge
    agent signatures, impersonating legitimate governance participants.
  \item \textbf{Historical decryption}: Messages encrypted with classical
    key exchange can be stored now and decrypted later when CRQCs become
    available (``harvest now, decrypt later'').
  \item \textbf{Credential forgery}: Consciousness credentials signed with
    classical keys become forgeable, undermining the entire governance tier
    system.
\end{enumerate}

NIST's post-quantum cryptography standardization~\cite{nist2024pqc} provides
the building blocks for quantum-resistant systems. We apply these to Mycelix's
governance infrastructure.

\section{Threat Model}

We consider an adversary with the following capabilities:

\begin{itemize}[nosep]
  \item Access to a future CRQC capable of running Shor's algorithm
  \item Ability to record all DHT traffic (``harvest now'')
  \item Participation in the Holochain network as a valid agent
  \item Cannot compromise a majority of DHT validators simultaneously
\end{itemize}

Our security goals:

\begin{enumerate}[nosep]
  \item \textbf{Confidentiality}: Messages remain private even against a
    future CRQC.
  \item \textbf{Integrity}: Source chain entries and credentials cannot be
    forged.
  \item \textbf{Forward secrecy}: Compromise of long-term keys does not
    reveal past communications.
  \item \textbf{Backward compatibility}: Agents without PQ keys can still
    participate (with reduced security guarantees).
\end{enumerate}

\section{Hybrid Key Infrastructure}

\subsection{DID Document Keys}

Each Mycelix DID document contains two key types:

\begin{description}[nosep]
  \item[Authentication] Ed25519 (32-byte public key). Used for source chain
    signing, DHT validation, and Holochain's native operations.
  \item[Key Agreement] ML-KEM-768 (1,184-byte public key). Used for encrypted
    communication, PQC mail, and future signature verification.
\end{description}

\subsection{ML-KEM-768 Properties}

ML-KEM-768 (FIPS 203)~\cite{nist2024pqc} is a lattice-based key encapsulation
mechanism:

\begin{table}[h]
\centering
\begin{tabular}{lr}
\toprule
\textbf{Parameter} & \textbf{Value} \\
\midrule
Security level & NIST Category 3 (AES-192) \\
Public key & 1,184 bytes \\
Secret key & 2,400 bytes \\
Ciphertext & 1,088 bytes \\
Shared secret & 32 bytes \\
Encapsulation & $<1$ ms \\
Decapsulation & $<1$ ms \\
\bottomrule
\end{tabular}
\caption{ML-KEM-768 parameters.}
\end{table}

\subsection{Key Size Impact on DHT}

The 37$\times$ increase in public key size (1,184 vs. 32 bytes) has
measurable but manageable impact on DHT operations:

\begin{itemize}[nosep]
  \item \textbf{Storage}: DID documents grow by $\sim$1.2 KB per agent.
    With 3$\times$ DHT replication, each agent's PQ key adds $\sim$3.6 KB
    to the network.
  \item \textbf{Bandwidth}: DID document retrieval (for key exchange) adds
    $\sim$1.2 KB per query. This is negligible compared to typical entry
    sizes (1--10 KB).
  \item \textbf{Gossip}: DHT gossip propagating DID updates transmits the
    full document. PQ keys add $\sim$1.2 KB per gossip message for DID
    updates specifically.
\end{itemize}

For a network of 10,000 agents, the total PQ key storage overhead is
$\sim$36 MB (with replication)---negligible for modern hardware.

\section{Hybrid Handshake Protocol}

\subsection{Protocol Description}

The hybrid handshake combines classical X25519 Diffie-Hellman with ML-KEM-768
key encapsulation:

\begin{algorithm}[h]
\caption{Hybrid Handshake}
\begin{algorithmic}[1]
\State \textbf{Initiator} $A$, \textbf{Responder} $B$
\State $A$: Generate ephemeral X25519 key pair $(ek_A, EK_A)$
\State $A$: Retrieve $B$'s ML-KEM-768 public key $PK_B^{\text{PQ}}$ from DHT
\State $A$: $(ct, ss_{\text{PQ}}) \gets \text{ML-KEM.Encaps}(PK_B^{\text{PQ}})$
\State $A \to B$: $\{EK_A, ct\}$
\State $B$: Generate ephemeral X25519 key pair $(ek_B, EK_B)$
\State $B$: $ss_{\text{classical}} \gets \text{X25519}(ek_B, EK_A)$
\State $B$: $ss_{\text{PQ}} \gets \text{ML-KEM.Decaps}(sk_B^{\text{PQ}}, ct)$
\State $B \to A$: $\{EK_B\}$
\State $A$: $ss_{\text{classical}} \gets \text{X25519}(ek_A, EK_B)$
\State \textbf{Both}: $K \gets \text{HKDF}(ss_{\text{classical}} \| ss_{\text{PQ}})$
\end{algorithmic}
\end{algorithm}

The final session key $K$ is derived from both shared secrets using HKDF.
This provides IND-CCA2 security if \emph{either} X25519 or ML-KEM-768 is
secure.

\begin{theorem}[Hybrid security]
If either X25519 (under the CDH assumption) or ML-KEM-768 (under the
Module-LWE assumption) is IND-CCA2 secure, then the hybrid handshake
produces a session key indistinguishable from random.
\end{theorem}

This follows from the standard hybrid argument: an adversary that can
distinguish the session key from random must break \emph{both} underlying
schemes.

\subsection{Forward Secrecy}

Forward secrecy is provided by the ephemeral X25519 key pairs. Even if
long-term ML-KEM-768 keys are compromised, past session keys cannot be
recovered because the ephemeral X25519 keys are erased after key derivation.

\subsection{Fallback for Non-PQ Agents}

Agents without ML-KEM-768 keys (not yet upgraded) fall back to classical-only
X25519 key exchange. The system maintains backward compatibility while
encouraging upgrade through:

\begin{itemize}[nosep]
  \item Visual indicators showing PQ status in the UI
  \item Consciousness profile bonus for PQ key registration
  \item Governance recommendations to require PQ keys for Steward+ tier
\end{itemize}

\section{PQC-Encrypted Mail}

\subsection{12-Zome Architecture}

The Mail cluster implements decentralized email with 12 zomes:

\begin{enumerate}[nosep]
  \item \textbf{compose}: Message composition and encryption
  \item \textbf{inbox}: Incoming message management
  \item \textbf{outbox}: Sent message tracking
  \item \textbf{folders}: Custom folder management
  \item \textbf{contacts}: Contact list with PQ key caching
  \item \textbf{attachments}: Encrypted attachment storage
  \item \textbf{threads}: Conversation threading
  \item \textbf{search}: Encrypted search over local messages
  \item \textbf{spam}: Consciousness-gated spam filtering
  \item \textbf{keys}: Key management and rotation
  \item \textbf{notifications}: Push notification handling
  \item \textbf{bridge}: Cross-cluster mail dispatch
\end{enumerate}

\subsection{Message Encryption}

Each message is encrypted using AES-256-GCM with a key derived from the
hybrid handshake:

\begin{enumerate}[nosep]
  \item Initiator performs hybrid handshake with recipient (Section 4)
  \item Session key $K$ is derived via HKDF
  \item Message plaintext is encrypted: $C = \text{AES-256-GCM}(K, \text{nonce}, M)$
  \item Encrypted message is stored on sender's source chain (outbox) and
    sent to recipient's DHT neighborhood (inbox)
  \item Recipient decapsulates their ML-KEM-768 ciphertext to recover $K$
    and decrypts the message
\end{enumerate}

\subsection{Conversation Forward Secrecy}

Each conversation uses a ratcheting key derivation:

\begin{equation}
  K_{n+1} = \text{HKDF}(K_n \| \text{message\_hash}_n)
\end{equation}

This ensures that compromise of key $K_n$ does not reveal messages encrypted
with $K_{n-1}$ or earlier, providing conversation-level forward secrecy.

\subsection{Spam Prevention}

Mail sending requires Participant+ consciousness tier. This is the most
effective decentralized spam prevention: creating a Participant-tier identity
requires non-trivial identity verification ($I \geq 0.25$) and community
engagement, making spam account creation expensive.

Additionally, the spam zome implements:
\begin{itemize}[nosep]
  \item Rate limiting per sender (configurable per recipient)
  \item Reputation-weighted spam scoring (low-reputation senders are flagged)
  \item Community blocklists (shared among trusted peers)
\end{itemize}

\section{Reputation System Security}

\subsection{Temporal Decay}

Reputation scores decay exponentially:

\begin{equation}
  R(t) = R_0 \cdot 0.998^{(t - t_0)/\text{day}}
\end{equation}

This prevents an agent from accumulating permanent reputation from past
good behavior and then exploiting it. The half-life of
$\frac{\ln 2}{\ln(1/0.998)} \approx 346$ days means reputation approximately
halves yearly.

\subsection{Slashing}

Severe misbehavior (detected by safety agents or community reports) triggers
reputation slashing:

\begin{equation}
  R \gets 0.5 \cdot R
\end{equation}

Multiple slashing events compound: three slashing events reduce reputation
to $0.125 \cdot R_0$, pushing most agents below the blacklist threshold
of 0.05.

\subsection{Blacklisting and Restoration}

Agents with $R < 0.05$ are blacklisted from all governance participation.
Restoration requires 100 verified positive interactions---approximately
3--4 months of active, positive community participation.

\section{Implementation Considerations}

\subsection{PQC Library Selection}

We use the \texttt{ml-kem} Rust crate (pure Rust implementation of
FIPS 203). Key considerations:

\begin{itemize}[nosep]
  \item \textbf{WASM compatibility}: The crate compiles to
    \texttt{wasm32-unknown-unknown} for Holochain zomes.
  \item \textbf{Constant-time operations}: The implementation uses
    constant-time comparisons to prevent timing side-channels.
  \item \textbf{No system dependencies}: Pure Rust avoids C library
    linking issues in the WASM target.
\end{itemize}

\subsection{Key Storage}

ML-KEM-768 secret keys (2,400 bytes) are stored in the agent's secure
keystore, managed by Holochain's conductor. The conductor's keystore
provides:

\begin{itemize}[nosep]
  \item Encrypted at-rest storage
  \item Memory protection (mlock where available)
  \item Key derivation from a master secret
\end{itemize}

\section{Evaluation}

\subsection{Performance}

ML-KEM-768 encapsulation/decapsulation adds $<2$ ms to message send/receive
operations. This is negligible compared to DHT propagation latency (typically
100--500 ms).

\subsection{Bandwidth Overhead}

Per-message overhead from PQC:
\begin{itemize}[nosep]
  \item ML-KEM-768 ciphertext: 1,088 bytes (first message only, then ratcheted)
  \item AES-256-GCM: 16-byte tag + 12-byte nonce = 28 bytes per message
  \item Total first-message overhead: $\sim$1.1 KB
  \item Subsequent message overhead: 28 bytes
\end{itemize}

\subsection{Migration Path}

The hybrid architecture supports incremental migration:

\begin{enumerate}[nosep]
  \item \textbf{Phase 1} (current): Classical keys required, PQ keys optional.
    Hybrid handshake used when both parties have PQ keys.
  \item \textbf{Phase 2} (planned): PQ keys recommended for Citizen+ tier.
    Governance actions may require PQ keys.
  \item \textbf{Phase 3} (future): PQ keys required for all agents. Classical
    keys retained for backward compatibility only.
\end{enumerate}

\section{Related Work}

\textbf{Signal Protocol}~\cite{signal}: Provides forward secrecy and
ratcheting but uses classical-only key exchange. The PQXDH
proposal~\cite{signal_pqxdh} adds post-quantum keys but is centralized
(Signal servers).

\textbf{Matrix/Olm}~\cite{matrix}: Federated encrypted messaging with
classical-only keys. No PQC migration plan published.

\textbf{ProtonMail}: Server-mediated PQC encryption. Not decentralized.

Our contribution is the first PQC-encrypted messaging system that is
(1) fully decentralized (Holochain DHT, no servers), (2) hybrid
(classical + PQ), and (3) integrated with consciousness-gated access control.

\section{Conclusion}

Post-quantum privacy for decentralized governance is both necessary and
practical. The hybrid handshake protocol provides quantum resistance today
without sacrificing interoperability with classical infrastructure. The
12-zome PQC-encrypted mail system demonstrates that full-featured decentralized
email is achievable with acceptable performance overhead ($<2$ ms per operation,
$\sim$1.1 KB first-message bandwidth). The integration with consciousness
gating (Participant+ for mail, Citizen+ recommended for PQ keys) provides
Sybil-resistant spam prevention while encouraging adoption of quantum-resistant
cryptography.

\begin{thebibliography}{99}
  \bibitem{shor1994} Shor, P. W. (1994). Algorithms for Quantum Computation:
    Discrete Logarithms and Factoring. \textit{FOCS}.
  \bibitem{nist2024pqc} NIST (2024). Module-Lattice-Based Key-Encapsulation
    Mechanism Standard. FIPS 203 (ML-KEM).
  \bibitem{signal} Marlinspike, M. and Perrin, T. (2016). The X3DH Key
    Agreement Protocol. Signal Specification.
  \bibitem{signal_pqxdh} Signal (2023). PQXDH: Post-Quantum Extended
    Diffie-Hellman. Signal Specification.
  \bibitem{matrix} Matrix.org Foundation (2020). Olm: An Implementation of
    the Double Ratchet. Matrix Specification.
  \bibitem{brock2018} Brock, A. and Harris-Braun, E. (2018). Holochain:
    Scalable Agent-Centric Distributed Computing.
\end{thebibliography}

\end{document}
